\chapter{The Operational Amplifier}\label{ch:b1:operational-amplifier}

A strain gauge glued to a bridge girder changes its resistance by a
thousandth when a truck passes; the \omterm{prop:b1:dc-circuits:wheatstone}{Wheatstone bridge} around it
(\cref{ch:b1:dc-circuits}) turns that into two millivolts. Two
millivolts drive nothing. Between the gauge and the alarm that will
close the bridge to traffic sits a chip the size of a fingernail, sold
for a few cents, which multiplies the \omterm{def:b1:wave-propagation:signal}{signal} by a thousand, \omterm{def:b1:filters-transfer-functions:H}{filters} it,
compares it with a threshold and lights a lamp --- four jobs, one
component, four different wirings. This chapter introduces the
\omterm{def:b1:operational-amplifier:opamp}{operational amplifier} through its ideal model, derives the handful of
circuits that every instrument is built from, and shows what happens
when the feedback is removed.

\section{The component and its ideal model}

\begin{definition}[Operational amplifier]\label{def:b1:operational-amplifier:opamp}
\index{operational amplifier}\index{saturation voltage}\index{open-loop gain}
An \emph{operational amplifier} is an integrated circuit with two
inputs --- the \emph{non-inverting} input at potential $V_+$ and the
\emph{inverting} input at $V_-$ --- and one output at potential $V_s$,
powered by two supplies $\pm V_{\mathrm{cc}}$ (typically
$\pm\qty{15}{V}$) that are usually left off the diagrams. It amplifies
the difference $\epsilon = V_+ - V_-$:
\[
V_s = \mu\,\epsilon \quad\text{as long as } \abs{V_s} < V_{\mathrm{sat}},
\]
with an \emph{open-loop gain} $\mu$ of order $10^5$ and a
\emph{saturation \omterm{def:b1:dc-circuits:voltage}{voltage}} $V_{\mathrm{sat}}$ slightly below
$V_{\mathrm{cc}}$. When $\mu\abs\epsilon$ would exceed $V_{\mathrm{sat}}$,
the output sticks at $+V_{\mathrm{sat}}$ or $-V_{\mathrm{sat}}$: the
\emph{saturated} regime.
\end{definition}

\begin{definition}[Ideal operational amplifier]\label{def:b1:operational-amplifier:ideal}
\index{ideal operational amplifier}\index{virtual short circuit}\index{linear regime}
The \emph{ideal} model assumes: (i) no current enters either input,
$i_+ = i_- = 0$ (infinite input \omterm{def:b1:sinusoidal-impedance:impedance}{impedance}); (ii) the output imposes its
\omterm{def:b1:dc-circuits:voltage}{voltage} whatever the current drawn (zero output \omterm{def:b1:sinusoidal-impedance:impedance}{impedance}); (iii)
infinite gain, $\mu \to \infty$. In the \emph{linear regime} ($\abs{V_s}
< V_{\mathrm{sat}}$) the third assumption forces
\[
\epsilon = V_+ - V_- = 0 :
\]
the two inputs sit at the same potential (a ``virtual short circuit'')
although no current flows between them.
\end{definition}

\begin{omfigure}
\begin{tikzpicture}[font=\small]
  \node[op amp] (oa) at (0,0) {};
  \draw (oa.-) -- ++(-0.8,0) node[left] {$V_-$};
  \draw (oa.+) -- ++(-0.8,0) node[left] {$V_+$};
  \draw (oa.out) -- ++(0.6,0) node[right] {$V_s$};
  \draw (oa.up) -- ++(0,0.3) node[above, font=\footnotesize] {$+V_{\mathrm{cc}}$};
  \draw (oa.down) -- ++(0,-0.3) node[below, font=\footnotesize] {$-V_{\mathrm{cc}}$};
  \begin{scope}[xshift=6.2cm]
  \begin{axis}[omaxis, axis y line=left, width=6.4cm, height=4.6cm, xlabel={$\epsilon$}, ylabel={$V_s$},
      xmin=-2.2, xmax=2.2, ymin=-1.5, ymax=1.5, xtick={-0.15,0.15}, xticklabels={,},
      ytick={-1,1}, yticklabels={$-V_{\mathrm{sat}}$,$+V_{\mathrm{sat}}$}, anchor=origin, at={(0,0)}]
    \addplot[omThm, thick, domain=-2.2:-0.15, samples=2] {-1};
    \addplot[omThm, thick, domain=-0.15:0.15, samples=2] {x/0.15};
    \addplot[omThm, thick, domain=0.15:2.2, samples=2] {1};
    \node[omThm, font=\footnotesize] at (axis cs:1.3,0.6) {saturated};
    \node[omThm, font=\footnotesize] at (axis cs:-1.2,-0.6) {saturated};
    \node[omThm, font=\footnotesize, right] at (axis cs:0.2,-0.55) {linear, slope $\mu$};
  \end{axis}
  \end{scope}
\end{tikzpicture}

{\small The \omterm{def:b1:operational-amplifier:opamp}{operational amplifier}: symbol with its two inputs, and its
transfer characteristic $V_s(\epsilon)$ --- a steep linear segment of
slope $\mu \sim 10^5$ (drawn far too shallow), then saturation at
$\pm V_{\mathrm{sat}}$. In the ideal model the linear segment is
vertical: $\epsilon = 0$.}
\end{omfigure}

\begin{proposition}[Negative feedback and the linear regime]\label{prop:b1:operational-amplifier:feedback}
\index{negative feedback}\index{positive feedback}
If the output is connected back to the \emph{inverting} input through a
network (\omterm{prop:b1:operational-amplifier:feedback}{negative feedback}), the circuit has a stable linear \omterm{met:b1:dc-circuits:loadline}{operating
point} at $\epsilon = 0$ and the ideal-amplifier rules apply. If the
output is fed back to the non-inverting input (\omterm{prop:b1:operational-amplifier:feedback}{positive feedback}), or
not fed back at all, the amplifier saturates: $V_s = \pm V_{\mathrm{sat}}$
according to the sign of $\epsilon$.
\end{proposition}

\begin{proof}[Partial proof]
With \omterm{prop:b1:operational-amplifier:feedback}{negative feedback}, suppose $V_s$ rises a little above its
equilibrium value; through the feedback network $V_-$ rises, $\epsilon$
falls, and the amplifier pulls $V_s$ back down: the perturbation is
corrected. With \omterm{prop:b1:operational-amplifier:feedback}{positive feedback} the same perturbation raises $V_+$,
increases $\epsilon$ and drives $V_s$ further away until it saturates.
The full argument (a first-order model of the amplifier, $\mu(j\omega)
= \mu_0/(1 + j\omega/\omega_0)$, and the stability of the resulting
differential equation) is the subject of the Year 2 volume; the rule
above is what one uses.
\end{proof}

\begin{remark}[Real amplifiers]\label{rem:b1:operational-amplifier:real}
\index{gain-bandwidth product}\index{slew rate}\index{input offset voltage}
Three departures from the ideal matter in practice. The \omterm{def:b1:operational-amplifier:opamp}{open-loop gain}
falls with frequency as $\mu_0/(1 + jf/f_0)$ with $f_0 \sim \qty{10}{Hz}$:
the product $\mu_0 f_0 = f_T \sim \qty{1}{MHz}$ (the \emph{gain--bandwidth
product}) is the bandwidth available to a follower, and a circuit of
closed-loop gain $G$ has bandwidth $f_T/G$. The output cannot change
faster than the \emph{\omterm{rem:b1:operational-amplifier:real}{slew rate}}, $\sim\qty{1}{V/\micro s}$. And a
small \emph{\omterm{rem:b1:operational-amplifier:real}{input offset voltage}} ($\sim\qty{1}{mV}$) is amplified like
a \omterm{def:b1:wave-propagation:signal}{signal} --- fatal for millivolt inputs unless trimmed.
\end{remark}

\section{The basic linear circuits}

\begin{method}[Analyzing an ideal op-amp circuit in the linear regime]\label{met:b1:operational-amplifier:method}
\begin{enumerate}
  \item Check that the feedback goes to the inverting input.
  \item Write $i_+ = i_- = 0$: the input nodes obey Kirchhoff's \omterm{thm:b1:dc-circuits:kirchhoff}{node
        law} with no current into the amplifier (dividers and Millman's
        theorem apply).
  \item Write $V_+ = V_-$ and solve for $V_s$.
  \item Verify afterwards that $\abs{V_s} < V_{\mathrm{sat}}$; if not,
        the amplifier is saturated and the result is $\pm V_{\mathrm{sat}}$.
\end{enumerate}
\end{method}

\begin{proposition}[Follower, non-inverting and inverting amplifiers]\label{prop:b1:operational-amplifier:basic}
\index{voltage follower}\index{non-inverting amplifier}\index{inverting amplifier}\index{virtual ground}
\begin{itemize}
  \item \emph{Follower}: output wired to $V_-$, input on $V_+$:
        $V_s = V_e$; infinite input \omterm{def:b1:sinusoidal-impedance:impedance}{impedance}, zero output \omterm{def:b1:sinusoidal-impedance:impedance}{impedance}.
  \item \emph{\omterm{prop:b1:operational-amplifier:basic}{Non-inverting amplifier}}: output to $V_-$ through a
        divider $R_2$ (feedback) and $R_1$ (to ground), input on $V_+$:
        \[
        V_s = \left(1 + \frac{R_2}{R_1}\right)V_e .
        \]
  \item \emph{\omterm{prop:b1:operational-amplifier:basic}{Inverting amplifier}}: $V_+$ grounded, input through
        $R_1$ to $V_-$, feedback $R_2$ from output to $V_-$:
        \[
        V_s = -\frac{R_2}{R_1}\,V_e ,
        \]
        with $V_-$ held at ground potential (a \emph{\omterm{prop:b1:operational-amplifier:basic}{virtual ground}})
        and an input \omterm{def:b1:sinusoidal-impedance:impedance}{impedance} $R_1$.
\end{itemize}
\end{proposition}

\begin{proof}
Follower: $V_- = V_s$ and $V_+ = V_e$; $\epsilon = 0$ gives $V_s = V_e$.
Non-inverting: no current into $V_-$, so the divider gives $V_- =
V_sR_1/(R_1 + R_2)$; equate to $V_+ = V_e$. Inverting: $V_- = V_+ = 0$;
\omterm{thm:b1:dc-circuits:kirchhoff}{node law} at $V_-$ with $i_- = 0$: $(V_e - 0)/R_1 + (V_s - 0)/R_2 = 0$.
The source delivers $V_e/R_1$: input \omterm{def:b1:sinusoidal-impedance:impedance}{impedance} $R_1$.
\end{proof}

\begin{omfigure}
\begin{tikzpicture}[font=\small]
  % follower (+ input up)
  \node[op amp, noinv input up] (a) at (0,0) {};
  \draw (a.+) -- ++(-0.6,0) node[left] {$V_e$};
  \draw (a.out) -- ++(0.5,0) node[right] {$V_s$};
  \draw (a.out) ++(0.25,0) -- ++(0,-1.3) -| ($(a.-)+(-0.4,0)$) -- (a.-);
  \node[below, font=\footnotesize] at (0,-1.7) {follower: $V_s = V_e$};
  % non-inverting (+ input up)
  \begin{scope}[xshift=5.6cm]
  \node[op amp, noinv input up] (b) at (0,0) {};
  \draw (b.+) -- ++(-0.8,0) node[left] {$V_e$};
  \draw (b.out) -- ++(0.6,0) node[right] {$V_s$};
  \coordinate (j) at ($(b.-)+(-0.5,0)$);
  \draw (b.-) -- (j);
  \coordinate (f) at ($(j)+(0,-1.0)$);
  \coordinate (o) at ($(b.out)+(0.3,0)$);
  \draw (o) -- (o |- f) to[R=$R_2$] (f) -- (j);
  \draw (f) to[R=$R_1$] ++(0,-1.6) node[ground]{};
  \node[below, font=\footnotesize] at (0,-3.7) {non-inverting: $V_s = (1 + R_2/R_1)V_e$};
  \end{scope}
  % inverting (- input up)
  \begin{scope}[xshift=11.3cm]
  \node[op amp] (c) at (0,0) {};
  \draw (c.+) -- ++(-0.4,0) node[ground]{};
  \coordinate (m) at ($(c.-)+(-0.5,0)$);
  \draw (c.-) -- (m) to[R=$R_1$] ++(-1.5,0) node[left] {$V_e$};
  \draw (c.out) -- ++(0.6,0) node[right] {$V_s$};
  \coordinate (f2) at ($(m)+(0,1.3)$);
  \coordinate (o2) at ($(c.out)+(0.3,0)$);
  \draw (o2) -- (o2 |- f2) to[R=$R_2$] (f2) -- (m);
  \node[below, font=\footnotesize] at (0,-1.7) {inverting: $V_s = -(R_2/R_1)V_e$};
  \end{scope}
\end{tikzpicture}

{\small The three basic amplifiers. In each the output feeds the
inverting input, the amplifier works in its \omterm{def:b1:operational-amplifier:ideal}{linear regime}, and the
gain is set by \omterm{def:b1:dc-circuits:resistor}{resistor} ratios alone --- not by the amplifier's own
enormous and ill-defined $\mu$.}
\end{omfigure}

\begin{example}[Why a follower]\label{ex:b1:operational-amplifier:follower}
A sensor of Thévenin resistance \qty{10}{k\ohm} feeding a \qty{1}{k\ohm}
load delivers only $1/11$ of its open-circuit \omterm{def:b1:dc-circuits:voltage}{voltage}
(\cref{rem:b1:dc-circuits:loading}). A follower between them draws no
current from the sensor and drives the load from a \omterm{def:b1:sinusoidal-impedance:impedance}{zero-impedance}
output: the full \omterm{def:b1:dc-circuits:voltage}{voltage} arrives. The follower is an \emph{\omterm{def:b1:sinusoidal-impedance:impedance}{impedance}
adapter} --- and it makes the cascade rule of
\cref{prop:b1:filters-transfer-functions:cascade} exact.
\end{example}

\begin{proposition}[Summing and difference amplifiers]\label{prop:b1:operational-amplifier:sumdiff}
\index{summing amplifier}\index{difference amplifier}
Inputs $V_1, V_2$ through $R_1, R_2$ to the inverting node, feedback
$R_f$, $V_+$ grounded:
\[
V_s = -R_f\left(\frac{V_1}{R_1} + \frac{V_2}{R_2}\right)
\quad (\text{summing amplifier}).
\]
With $V_1$ through $R_1$ to $V_-$ (feedback $R_2$) and $V_2$ through
$R_1$ to $V_+$ (with $R_2$ from $V_+$ to ground):
\[
V_s = \frac{R_2}{R_1}\,(V_2 - V_1)
\quad (\text{difference amplifier}).
\]
\end{proposition}

\begin{proof}
Summing: \omterm{thm:b1:dc-circuits:kirchhoff}{node law} at the \omterm{prop:b1:operational-amplifier:basic}{virtual ground}, $V_1/R_1 + V_2/R_2 + V_s/R_f =
0$. Difference: $V_+ = V_2R_2/(R_1 + R_2)$ (divider); \omterm{thm:b1:dc-circuits:kirchhoff}{node law} at $V_-$:
$(V_1 - V_-)/R_1 + (V_s - V_-)/R_2 = 0$, so $V_s = V_-(1 + R_2/R_1) -
V_1R_2/R_1$; set $V_- = V_+$ and simplify.
\end{proof}

\section{Integrator and differentiator}

\begin{proposition}[Integrator]\label{prop:b1:operational-amplifier:integrator}
\index{integrator (op-amp)}
Input through $R$ to $V_-$, a capacitor $C$ as feedback, $V_+$
grounded:
\[
V_s(t) = -\frac{1}{RC}\int_0^t V_e(t')\,\dd t' + V_s(0),
\qquad
\underline H = -\frac{1}{jRC\omega} .
\]
The gain falls \qty{20}{dB} per decade at every frequency and the phase
is $+90^\circ$: an exact integrator --- until the smallest offset or
bias, integrated forever, drives the output into saturation.
\end{proposition}

\begin{proof}
\omterm{prop:b1:operational-amplifier:basic}{Virtual ground}: the current $V_e/R$ flows into the capacitor, whose
\omterm{def:b1:dc-circuits:voltage}{voltage} (from $V_-$ to the output) is $-V_s$; hence $C\,\dd(-V_s)/\dd t
= V_e/R$, and in complex notation $\underline H = -1/(jRC\omega)$.
\end{proof}

\begin{remark}[Taming the integrator]\label{rem:b1:operational-amplifier:pseudo}
\index{pseudo-integrator}
A \omterm{def:b1:dc-circuits:resistor}{resistor} $R'$ in parallel with $C$ gives
\[
\underline H = -\frac{R'/R}{1 + jR'C\omega},
\]
a first-order low-pass of DC gain $-R'/R$ and cutoff $1/R'C$, which
integrates for $\omega \gg 1/R'C$ but cannot drift to saturation. Likewise the differentiator (capacitor at the input,
\omterm{def:b1:dc-circuits:resistor}{resistor} as feedback: $V_s = -RC\,\dd V_e/\dd t$, $\underline H =
-jRC\omega$) amplifies high-frequency noise without limit unless a
small series \omterm{def:b1:dc-circuits:resistor}{resistor} caps its gain.
\end{remark}

\begin{omfigure}
\begin{tikzpicture}[font=\small]
  \node[op amp] (c) at (0,0) {};
  \draw (c.+) -- ++(-0.4,0) node[ground]{};
  \draw (c.-) -- ++(-0.3,0) coordinate (m1) to[R=$R$] ++(-1.8,0) node[left] {$V_e$};
  \draw (c.out) -- ++(0.5,0) node[right] {$V_s$};
  \draw (c.out) ++(0.25,0) |- ++(0,1.2) to[C=$C$] ++(-2.1,0) -| (m1);
  \begin{scope}[xshift=7.3cm, yshift=-1.2cm]
  \begin{axis}[omaxis, axis x line=bottom, axis y line=left, xlabel style={at={(axis description cs:0.5,-0.16)}, anchor=north},
      width=6.4cm, height=4.2cm, xmode=log, xlabel={$\omega$ ($\log$)}, ylabel={$G_{\mathrm{dB}}$},
      xmin=0.01, xmax=100, ymin=-45, ymax=48, ytick={40,20,0,-20,-40}, xtick={0.01,0.1,1,10,100},
      xticklabels={$0.01$,$0.1$,$1$,$10$,$100$}, anchor=origin, at={(0,0)}]
    \addplot[omDef, thick, domain=0.01:100, samples=2] {-20*log10(x)};
    \addplot[omThm, thick, domain=0.01:100, samples=200] {20 - 10*log10(1 + (x/0.1)^2)};
    \node[omDef, font=\footnotesize, rotate=-32] at (axis cs:8,-10) {integrator};
    \node[omThm, font=\footnotesize, anchor=west] at (axis cs:0.012,29) {pseudo-integrator};
  \end{axis}
  \end{scope}
\end{tikzpicture}

{\small The integrator and its \omterm{def:b1:filters-transfer-functions:bode}{Bode diagram} (blue, \qty{-20}{dB/decade}
throughout, infinite gain at DC); with a \omterm{def:b1:dc-circuits:resistor}{resistor} across $C$ the gain
is capped at low frequency (red): a low-pass that integrates above its
cutoff.}
\end{omfigure}

\begin{example}[Triangle from square]\label{ex:b1:operational-amplifier:triangle}
$R = \qty{10}{k\ohm}$, $C = \qty{100}{nF}$, $1/RC = \qty{1000}{s^{-1}}$.
A \qty{1}{kHz} square wave of $\pm\qty{1}{V}$ gives a triangle of slope
$\mp\qty{1000}{V/s}$ and \omterm{def:b1:wave-propagation:signal}{amplitude} $ET/4RC = \qty{0.25}{V}$ --- without
the $1/\omega$ attenuation of a passive RC, and with a \omterm{def:b1:sinusoidal-impedance:impedance}{zero-impedance}
output to feed the next stage.
\end{example}

\section{Active filters}

\begin{proposition}[First-order active low-pass]\label{prop:b1:operational-amplifier:activelp}
\index{active filter}
The \omterm{prop:b1:operational-amplifier:basic}{inverting amplifier} with a capacitor $C$ across its feedback
\omterm{def:b1:dc-circuits:resistor}{resistor} $R_2$ has
\[
\underline H = -\frac{R_2/R_1}{1 + jR_2C\omega}:
\]
a low-pass of DC gain $-R_2/R_1$ (which can exceed $1$) and cutoff
$\omega_c = 1/R_2C$, with zero output \omterm{def:b1:sinusoidal-impedance:impedance}{impedance} --- stages cascade by
simple multiplication.
\end{proposition}

\begin{proof}
\omterm{prop:b1:operational-amplifier:basic}{Inverting amplifier} with the feedback \omterm{def:b1:sinusoidal-impedance:impedance}{impedance} $R_2 \parallel 1/jC\omega
= R_2/(1 + jR_2C\omega)$ in place of $R_2$.
\end{proof}

\section{Without negative feedback: comparators}

\begin{proposition}[Comparator]\label{prop:b1:operational-amplifier:comparator}
\index{comparator}
With no feedback, the ideal amplifier is saturated:
$V_s = +V_{\mathrm{sat}}$ if $V_+ > V_-$, $-V_{\mathrm{sat}}$ if
$V_+ < V_-$. With a reference $V_{\mathrm{ref}}$ on one input and a
\omterm{def:b1:wave-propagation:signal}{signal} on the other, the output tells whether the \omterm{def:b1:wave-propagation:signal}{signal} exceeds the
reference --- a one-bit converter.
\end{proposition}

\begin{proof}
$\mu\epsilon$ exceeds $V_{\mathrm{sat}}$ for any $\abs\epsilon >
V_{\mathrm{sat}}/\mu \sim \qty{0.1}{mV}$.
\end{proof}

\begin{proposition}[Schmitt trigger]\label{prop:b1:operational-amplifier:schmitt}
\index{Schmitt trigger}\index{hysteresis}
Feed a fraction $\beta = R_1/(R_1 + R_2)$ of the output back to $V_+$
(divider $R_1$ to ground, $R_2$ from the output) and apply the \omterm{def:b1:wave-propagation:signal}{signal}
$V_e$ to $V_-$. The output switches from $+V_{\mathrm{sat}}$ to
$-V_{\mathrm{sat}}$ when $V_e$ rises past $+\beta V_{\mathrm{sat}}$, and
back when $V_e$ falls below $-\beta V_{\mathrm{sat}}$: two thresholds,
a \emph{hysteresis} of width $2\beta V_{\mathrm{sat}}$, immune to noise
smaller than that width.
\end{proposition}

\begin{proof}
\omterm{prop:b1:operational-amplifier:feedback}{Positive feedback}: saturated. If $V_s = +V_{\mathrm{sat}}$ then $V_+ =
\beta V_{\mathrm{sat}}$, and $\epsilon = \beta V_{\mathrm{sat}} - V_e$
stays positive until $V_e$ exceeds $\beta V_{\mathrm{sat}}$; then the
output flips, $V_+$ drops to $-\beta V_{\mathrm{sat}}$, and the new state
holds until $V_e$ falls below that.
\end{proof}

\begin{omfigure}
\begin{tikzpicture}[font=\small]
  \node[op amp] (s) at (0,0) {};
  \draw (s.-) -- ++(-0.8,0) node[left] {$V_e$};
  \draw (s.out) -- ++(0.6,0) node[right] {$V_s$};
  \coordinate (j) at ($(s.+)+(-0.6,0)$);
  \coordinate (f) at ($(j)+(0,-1.0)$);
  \coordinate (o) at ($(s.out)+(0.3,0)$);
  \draw (s.+) -- (j);
  \draw (o) -- (o |- f) to[R=$R_2$] (f) -- (j);
  \draw (f) to[R=$R_1$] ++(0,-1.6) node[ground]{};
  \begin{scope}[xshift=7.8cm, yshift=-0.9cm]
  \begin{axis}[omaxis, axis x line=bottom, axis y line=left, width=6.4cm, height=4.6cm, xlabel={$V_e$}, ylabel={$V_s$},
      xlabel style={at={(axis description cs:0.5,-0.12)}, anchor=north},
      xmin=-2.2, xmax=2.2, ymin=-1.5, ymax=1.5, xtick={-1,1}, xticklabels={$-\beta V_{\mathrm{sat}}$,$\beta V_{\mathrm{sat}}$},
      ytick={-1,1}, yticklabels={$-V_{\mathrm{sat}}$,$+V_{\mathrm{sat}}$}, anchor=origin, at={(0,0)}]
    \draw[omThm, thick, -Stealth] (axis cs:-2.2,1) -- (axis cs:0.2,1);
    \draw[omThm, thick] (axis cs:0.2,1) -- (axis cs:1,1) -- (axis cs:1,-1);
    \draw[omThm, thick, -Stealth] (axis cs:2.2,-1) -- (axis cs:-0.2,-1);
    \draw[omThm, thick] (axis cs:-0.2,-1) -- (axis cs:-1,-1) -- (axis cs:-1,1);
  \end{axis}
  \end{scope}
\end{tikzpicture}

{\small The \omterm{prop:b1:operational-amplifier:schmitt}{Schmitt trigger}: \omterm{prop:b1:operational-amplifier:feedback}{positive feedback} through $R_1$, $R_2$
gives two thresholds $\pm\beta V_{\mathrm{sat}}$. The output flips
high-to-low when $V_e$ rises past the upper threshold and low-to-high
when it falls below the lower one (arrows): a hysteresis cycle.}
\end{omfigure}

\begin{example}[The relaxation oscillator]\label{ex:b1:operational-amplifier:astable}
\index{relaxation oscillator}
Connect the \omterm{prop:b1:operational-amplifier:schmitt}{Schmitt trigger}'s output back to its own inverting input
through $R$, with $C$ from $V_-$ to ground. The capacitor charges
toward $\pm V_{\mathrm{sat}}$ and flips the trigger each time it reaches
$\pm\beta V_{\mathrm{sat}}$: the output is a square wave, the capacitor
\omterm{def:b1:dc-circuits:voltage}{voltage} a chain of exponential arcs, and the period is
\[
T = 2RC\ln\frac{1 + \beta}{1 - \beta}
\]
(\cref{exo:b1:operational-amplifier:12}): $2.2\,RC$ for $\beta = \tfrac12$.
No input, one capacitor, and a clock --- a system that oscillates
because it cannot settle, a theme that returns with the feedback
oscillators of the Year 2 volume.
\end{example}

\section{Exercises}

\begin{exercise}[$\star$]\label{exo:b1:operational-amplifier:1}
A \omterm{prop:b1:operational-amplifier:basic}{non-inverting amplifier} has $R_1 = \qty{1.0}{k\ohm}$, $R_2 =
\qty{9.0}{k\ohm}$, $V_{\mathrm{sat}} = \qty{15}{V}$. Gain? Output for
$V_e = \qty{0.50}{V}$? Largest input before saturation?
\end{exercise}

\begin{exercise}[$\star$]\label{exo:b1:operational-amplifier:2}
An \omterm{prop:b1:operational-amplifier:basic}{inverting amplifier} has $R_1 = \qty{10}{k\ohm}$, $R_2 = \qty{100}{k\ohm}$.
Gain, input \omterm{def:b1:sinusoidal-impedance:impedance}{impedance}, current drawn from a \qty{1.0}{V} source, output.
\end{exercise}

\begin{exercise}[$\star$]\label{exo:b1:operational-amplifier:3}
A sensor of \omterm{def:b1:dc-circuits:sources}{internal resistance} \qty{10}{k\ohm} and open-circuit
\omterm{def:b1:dc-circuits:voltage}{voltage} \qty{1.0}{V} must drive a \qty{1.0}{k\ohm} load. \omterm{def:b1:dc-circuits:voltage}{Voltage} across
the load without and with a follower between them?
\end{exercise}

\begin{exercise}[$\star$]\label{exo:b1:operational-amplifier:4}
A \omterm{prop:b1:operational-amplifier:sumdiff}{summing amplifier} has $R_f = \qty{10}{k\ohm}$ and inputs through
\qty{10}{k\ohm} and \qty{20}{k\ohm}. Write $V_s$; compute it for
$V_1 = \qty{1.0}{V}$, $V_2 = \qty{2.0}{V}$.
\end{exercise}

\begin{exercise}[$\star\star$]\label{exo:b1:operational-amplifier:5}
Derive the \omterm{prop:b1:operational-amplifier:sumdiff}{difference amplifier}'s $V_s = (R_2/R_1)(V_2 - V_1)$. With
$R_1 = \qty{10}{k\ohm}$, $R_2 = \qty{100}{k\ohm}$, $V_1 = \qty{1.00}{V}$,
$V_2 = \qty{1.05}{V}$: output? What happened to the \qty{1}{V} common
to both inputs?
\end{exercise}

\begin{exercise}[$\star\star$]\label{exo:b1:operational-amplifier:6}
Integrator with $R = \qty{10}{k\ohm}$, $C = \qty{100}{nF}$, output
initially zero. A \qty{1.0}{V} step is applied: slope of the output and
time to saturation (\qty{15}{V}). A \qty{1.0}{kHz} square wave of
$\pm\qty{1.0}{V}$: shape and \omterm{def:b1:wave-propagation:signal}{amplitude} of the output.
\end{exercise}

\begin{exercise}[$\star\star$]\label{exo:b1:operational-amplifier:7}
A \qty{1.0}{M\ohm} \omterm{def:b1:dc-circuits:resistor}{resistor} is added across the capacitor of the previous
integrator. Give the \omterm{def:b1:filters-transfer-functions:H}{transfer function}, the DC gain, the \omterm{def:b1:filters-transfer-functions:bode}{cutoff
frequency}, and the gain at \qty{1.0}{kHz}; in which range does the
circuit still integrate?
\end{exercise}

\begin{exercise}[$\star\star$]\label{exo:b1:operational-amplifier:8}
Design an inverting active low-pass with a gain of \qty{20}{dB} in the
\omterm{def:b1:filters-transfer-functions:bode}{passband} and a cutoff at \qty{500}{Hz}, using $R_1 = \qty{1.0}{k\ohm}$.
\end{exercise}

\begin{exercise}[$\star\star$]\label{exo:b1:operational-amplifier:9}
A real amplifier has $\mu_0 = 10^5$ and $f_T = \qty{1.0}{MHz}$. For a
\omterm{prop:b1:operational-amplifier:basic}{non-inverting amplifier} designed for a gain of $100$: actual DC gain
(keep $\mu_0$ finite in the analysis), and bandwidth. Bandwidth of a
follower?
\end{exercise}

\begin{exercise}[$\star\star\star$]\label{exo:b1:operational-amplifier:10}
A comparator with $V_{\mathrm{ref}} = \qty{2.0}{V}$ on $V_-$ watches a
slowly rising \omterm{def:b1:wave-propagation:signal}{signal} carrying \qty{50}{mV} of noise. Describe the output
as the \omterm{def:b1:wave-propagation:signal}{signal} crosses \qty{2.0}{V}. Remedy?
\end{exercise}

\begin{exercise}[$\star\star\star$]\label{exo:b1:operational-amplifier:11}
\omterm{prop:b1:operational-amplifier:schmitt}{Schmitt trigger} with $R_1 = \qty{1.0}{k\ohm}$, $R_2 = \qty{9.0}{k\ohm}$,
$V_{\mathrm{sat}} = \qty{15}{V}$: derive and compute the two thresholds.
Sketch $V_s$ for a triangular $V_e$ of \omterm{def:b1:wave-propagation:signal}{amplitude} \qty{3}{V}.
\end{exercise}

\begin{exercise}[$\star\star\star$]\label{exo:b1:operational-amplifier:12}
\omterm{ex:b1:operational-amplifier:astable}{Relaxation oscillator} with $R_1 = R_2$ ($\beta = \tfrac12$). Derive the
period $T = 2RC\ln\frac{1 + \beta}{1 - \beta}$ from the charging of $C$
through $R$ between the two thresholds; choose $R$ for \qty{1.0}{kHz}
with $C = \qty{100}{nF}$; sketch $V_s(t)$ and $V_C(t)$.
\end{exercise}

\section{Problem: An instrumentation chain for a strain gauge}

\begin{problem}[{Weekend problem --- from two millivolts on a bridge
girder to a red lamp: followers, a difference amplifier, an active
filter and a Schmitt trigger, and the two imperfections that decide
whether the chain works}]
\label{pb:b1:operational-amplifier:1}
A strain gauge of resistance $R(1 + \epsilon)$, $R = \qty{1.00}{k\ohm}$,
$0 \leq \epsilon \leq \num{2.0e-3}$, forms a \omterm{prop:b1:dc-circuits:wheatstone}{Wheatstone bridge} with
three fixed \qty{1.00}{k\ohm} \omterm{def:b1:dc-circuits:resistor}{resistors}, fed by $E = \qty{5.00}{V}$.
The bridge output $u = V_B - V_A$ (between the two midpoints) is
$u \approx E\epsilon/4$ (the sign is chosen by wiring). Amplifiers:
ideal unless stated; $V_{\mathrm{sat}} = \qty{15}{V}$; where needed,
$f_T = \qty{1.0}{MHz}$, \omterm{rem:b1:operational-amplifier:real}{slew rate} \qty{0.5}{V/\micro s}, input offset
\qty{1}{mV}.

\medskip
\noindent\textbf{Part I --- The bridge and its loading.}
\begin{enumerate}
  \item Recall why $u \approx E\epsilon/4$ for small $\epsilon$.
  \item Compute $u$ at full scale ($\epsilon = \num{2.0e-3}$).
  \item Give the potential of each midpoint (relative to the bridge's
        negative terminal) at zero strain and at full scale. What is
        their common value called?
  \item Each midpoint is the output of a divider of two \qty{1}{k\ohm}
        \omterm{def:b1:dc-circuits:resistor}{resistors}: what is the Thévenin resistance seen between the
        two midpoints?
  \item An amplifier whose input resistance is \qty{10}{k\ohm} is
        connected across the bridge output: what fraction of $u$ does it
        receive?
  \item Why does a follower on each midpoint cure this? State the two
        properties used.
  \item What are the followers' output \omterm{def:b1:dc-circuits:voltage}{voltages} (in terms of $V_A$,
        $V_B$), and what current do they draw from the bridge?
\end{enumerate}

\medskip
\noindent\textbf{Part II --- Amplifying the difference.}
\begin{enumerate}[resume]
  \item The followers feed a \omterm{prop:b1:operational-amplifier:sumdiff}{difference amplifier} ($R_1$ at both
        inputs, $R_2$ as feedback and to ground). Derive $V_s =
        (R_2/R_1)(V_B - V_A)$.
  \item Choose $R_1 = \qty{10}{k\ohm}$ and $R_2$ for a gain of $100$.
  \item A second, non-inverting stage multiplies by $10$: give its
        \omterm{def:b1:dc-circuits:resistor}{resistors}, the total gain, and the output at full scale.
  \item Both midpoints sit near $E/2 = \qty{2.5}{V}$ (the common-mode
        \omterm{def:b1:dc-circuits:voltage}{voltage}). What does an exactly matched \omterm{prop:b1:operational-amplifier:sumdiff}{difference amplifier} do
        with it? If one \omterm{def:b1:dc-circuits:resistor}{resistor} ratio is off by $1\%$, estimate the
        spurious input-referred \omterm{def:b1:dc-circuits:voltage}{voltage} and compare with the \omterm{def:b1:wave-propagation:signal}{signal}.
  \item The first amplifier has a \qty{1}{mV} input offset. What output
        error does it produce after the total gain? Conclude.
  \item With $f_T = \qty{1}{MHz}$, what is the bandwidth of each stage?
        Is it sufficient for a \omterm{def:b1:wave-propagation:signal}{signal} below \qty{5}{Hz}?
  \item The output must swing by \qty{2.5}{V} when a truck arrives in
        about \qty{10}{ms}: does the \omterm{rem:b1:operational-amplifier:real}{slew rate} limit?
\end{enumerate}

\medskip
\noindent\textbf{Part III --- Filtering.}
The amplified \omterm{def:b1:wave-propagation:signal}{signal} carries a \qty{50}{Hz} hum picked up from the
mains.
\begin{enumerate}[resume]
  \item Design an inverting active low-pass with gain $-1$ and cutoff
        \qty{10}{Hz}, using $R_1 = R_2 = \qty{16}{k\ohm}$: find $C$.
  \item Compute its attenuation at \qty{50}{Hz} in \omterm{def:b1:filters-transfer-functions:bode}{decibels}.
  \item Compute the phase shift at \qty{5}{Hz} and the corresponding
        delay. Acceptable for an alarm?
  \item Why is an \omterm{prop:b1:operational-amplifier:activelp}{active filter} preferable here to a passive RC
        (two reasons)?
  \item A second identical stage is cascaded: attenuation at \qty{50}{Hz}?
        Why is the product rule exact here?
\end{enumerate}

\medskip
\noindent\textbf{Part IV --- Threshold and alarm.}
The alarm must trip when the strain exceeds \num{1.6e-3}.
\begin{enumerate}[resume]
  \item To what output \omterm{def:b1:dc-circuits:voltage}{voltage} of the chain does this correspond? A
        comparator with this reference: output states.
  \item The filtered \omterm{def:b1:wave-propagation:signal}{signal} still carries \qty{50}{mV} of noise near
        the threshold. What does the plain comparator do?
  \item A \omterm{prop:b1:operational-amplifier:schmitt}{Schmitt trigger} with the reference on one input and feedback
        $\beta$: show that the thresholds are $V_{\mathrm{ref}} \pm
        \beta V_{\mathrm{sat}}$; choose $\beta$ for $\pm\qty{0.10}{V}$
        and give $R_1$, $R_2$.
  \item State the two thresholds in volts and in strain.
  \item The output drives an LED (\qty{2.0}{V}, \qty{13}{mA}) through a
        \omterm{def:b1:dc-circuits:resistor}{resistor}: compute it. What protects the LED when the output is
        at $-V_{\mathrm{sat}}$?
  \item Summarize the chain, stage by stage, and name the two
        non-idealities that set its real limits.
\end{enumerate}
\end{problem}
